% ============================================================
% presentaciones/avance1/slides.tex
% Primer avance — Proyecto ASE
% Compilar: pdflatex slides.tex
% ============================================================

\documentclass[aspectratio=169, 12pt]{beamer}

% ---------- Tema y colores IBERO ----------
\usetheme{Madrid}
\usecolortheme{default}
\definecolor{iberoazul}{RGB}{0, 56, 101}
\definecolor{iberoverde}{RGB}{0, 148, 115}
\setbeamercolor{structure}{fg=iberoazul}
\setbeamercolor{frametitle}{bg=iberoazul, fg=white}
\setbeamercolor{title}{fg=white, bg=iberoazul}
\setbeamercolor{block title}{bg=iberoverde, fg=white}

% ---------- Paquetes ----------
\usepackage[utf8]{inputenc}
\usepackage[spanish]{babel}
\usepackage{amsmath}
\usepackage{graphicx}
\usepackage{booktabs}

% ---------- Metadatos ----------
\title[Siniestralidad GMM]{Predicción de Siniestralidad en\\
Seguros de Gastos Médicos Mayores}
\subtitle{Avance 1: Planteamiento y EDA}
\author[García · Mendoza · Torres]{Ana García \and Carlos Mendoza \and Sofía Torres}
\institute[IBERO]{Universidad Iberoamericana Ciudad de México\\
Licenciatura en Actuaría}
\date{Primavera 2025}

% ============================================================
\begin{document}
% ============================================================

\begin{frame}
  \titlepage
\end{frame}

% ---- Tabla de contenidos -----------------------------------
\begin{frame}{Contenido}
  \tableofcontents
\end{frame}

% ---- Sección 1 ---------------------------------------------
\section{Motivación y Pregunta de Investigación}

\begin{frame}{Contexto: Mercado GMM en México}
  \begin{columns}
    \column{0.55\textwidth}
      \begin{itemize}
        \item El ramo GMM representa \textbf{28\% de primas totales} del sector vida y accidentes (CNSF, 2022)
        \item Crecimiento promedio anual: \textbf{8.3\%} en los últimos 5 años
        \item Desafío principal: inflación médica supera inflación general en \textbf{3–5 pp} anuales
        \item Siniestralidad promedio del sector: \textbf{68\%} (2022)
      \end{itemize}
    \column{0.40\textwidth}
      % \includegraphics[width=\textwidth]{figuras/primas_gmm_tendencia.pdf}
      \textit{[Figura: evolución de primas GMM]}
  \end{columns}
\end{frame}

\begin{frame}{Pregunta de Investigación}
  \begin{block}{Pregunta central}
    ¿Pueden los modelos GLM y gradient boosting mejorar la estimación de la \textbf{prima pura} respecto a los métodos actuariales tradicionales para carteras de GMM en México?
  \end{block}

  \vspace{0.5em}
  \textbf{Sub-preguntas:}
  \begin{enumerate}
    \item ¿Qué variables de riesgo tienen mayor poder predictivo sobre la frecuencia de siniestros?
    \item ¿La distribución Gamma es adecuada para modelar la severidad en esta cartera?
    \item ¿El modelo compuesto XGBoost supera estadísticamente al GLM Poisson-Gamma?
  \end{enumerate}
\end{frame}

% ---- Sección 2 ---------------------------------------------
\section{Descripción de los Datos}

\begin{frame}{Dataset: Cartera GMM 2019–2023}
  \begin{table}
    \centering
    \caption{Resumen de la cartera analizada}
    \begin{tabular}{lrr}
      \toprule
      \textbf{Característica} & \textbf{Pólizas} & \textbf{Siniestros} \\
      \midrule
      Total de registros & 14,823 & 41,670 \\
      Periodo & 2019–2023 & 2019–2023 \\
      \% sin siniestros & 61.4\% & — \\
      Monto promedio & — & \$42,300 MXN \\
      Suma asegurada mediana & \$3'000,000 MXN & — \\
      \bottomrule
    \end{tabular}
  \end{table}
\end{frame}

\begin{frame}{Variables de riesgo disponibles}
  \begin{columns}
    \column{0.48\textwidth}
      \textbf{Variables del contratante:}
      \begin{itemize}
        \item Edad (18–85 años)
        \item Sexo (M/F)
        \item Región (CDMX, Norte, Centro, Sur)
        \item Tipo de plan (Individual/Familiar)
      \end{itemize}
    \column{0.48\textwidth}
      \textbf{Variables del contrato:}
      \begin{itemize}
        \item Suma asegurada
        \item Deducible
        \item Coaseguro (\%)
        \item Antigüedad en cartera
        \item Nivel hospitalario (1–3)
      \end{itemize}
  \end{columns}
\end{frame}

% ---- Sección 3 ---------------------------------------------
\section{Análisis Exploratorio}

\begin{frame}{Distribución de frecuencia de siniestros}
  \begin{columns}
    \column{0.5\textwidth}
      % \includegraphics[width=\textwidth]{figuras/hist_frecuencia.pdf}
      \textit{[Figura: histograma de N]}

      \vspace{0.5em}
      \small Prueba de sobredispersión: $\hat{\phi} = 1.23$\\
      (leve sobredispersión respecto a Poisson)
    \column{0.5\textwidth}
      \begin{table}
        \small
        \begin{tabular}{lr}
          \toprule
          Estadístico & Valor \\
          \midrule
          Media $\bar{N}$ & 2.81 \\
          Varianza & 3.45 \\
          Mediana & 2 \\
          Máximo & 18 \\
          \% $N=0$ & 61.4\% \\
          \bottomrule
        \end{tabular}
      \end{table}
  \end{columns}
\end{frame}

\begin{frame}{Plan de trabajo}
  \begin{table}
    \centering
    \begin{tabular}{clc}
      \toprule
      \textbf{Semana} & \textbf{Actividad} & \textbf{Estado} \\
      \midrule
      1–2 & Obtención y limpieza de datos & {\color{green}✓} \\
      3–4 & Análisis exploratorio & {\color{green}✓} \\
      5–6 & Modelado de frecuencia (GLM + XGB) & {\color{orange}En curso} \\
      7–8 & Modelado de severidad & Pendiente \\
      9–10 & Modelo compuesto y tarificación & Pendiente \\
      11–12 & Escritura y presentación final & Pendiente \\
      \bottomrule
    \end{tabular}
  \end{table}
\end{frame}

% ---- Cierre ------------------------------------------------
\begin{frame}
  \centering
  \vspace{1em}
  {\Large \textbf{Gracias}}\\[1em]
  {\large Preguntas y comentarios}\\[2em]
  \small
  \texttt{ana.garcia@ibero.mx}\\
  \texttt{carlos.mendoza@ibero.mx}\\
  \texttt{sofia.torres@ibero.mx}\\[1em]
  Código y datos: \texttt{github.com/usuario/ase-siniestralidad-gmm}
\end{frame}

\end{document}
